% !TeX program = xelatex
\documentclass[margin = 5mm, tikz]{standalone}
\usepackage{xfrac}
\usepackage[MnSymbol]{mathspec}
\usepackage{esvect}
\setmainfont[Mapping=tex-text,Ligatures={NoRequired,NoCommon,NoContextual}]{Comic Sans MS}
\setmathfont(Digits,Latin,Greek)[Numbers={Lining,Proportional}]{Comic Sans MS}
\usetikzlibrary{shapes.geometric,decorations,decorations.pathmorphing}
\usetikzlibrary{intersections}
\usetikzlibrary{patterns}
\usetikzlibrary{calc}
\usetikzlibrary{shapes.misc}
\usetikzlibrary{spy}

\pgfdeclarelayer{bg}
\pgfdeclarelayer{bbg}
\pgfsetlayers{bbg, bg, main}

\definecolor{c1}{HTML}{ac0000}
\definecolor{c2}{HTML}{00a0d5}

\newcommand{\abs}[1]{{\left|{#1}\right|}}

\tikzset{
	point/.style n args = {1}{circle, draw = black, inner sep = #1pt, fill = black},
	empty point/.style = {circle, draw = black, inner sep = 2pt, fill = white},
	cross/.style = {cross out, draw, minimum size = 2 * (#1 - \pgflinewidth), inner sep = 0pt, outer sep = 0pt},
	xkcd/.style n args = {1}{decorate, decoration = {random steps, pre length = #1mm, post length = #1mm, segment length = 0.6mm, amplitude = 0.2pt}},
	xkcd/.default = {4},
	point/.default = {2}
}

\begin{document}
	\begin{tikzpicture}
		\pgfmathsetmacro{\m}{6/sqrt(2)}
		\pgfmathsetmacro{\d}{6 - 6/sqrt(2)}
		\draw[xkcd, thick] (0,0) -- (6,0) -- (6,6) -- (0,6) -- cycle;
		\draw[xkcd, thick, fill = c1, fill opacity = 0.1] (0,0) -- (\m,0) -- (\m,\m) -- (0,\m) -- cycle;
		\draw[xkcd, thick, fill = c1, fill opacity = 0.1, xshift = \d cm, yshift = \d cm] (0,0) -- (\m,0) -- (\m,\m) -- (0,\m) -- cycle;
		\node[above] at (\d/2, 6) {$m-n$};
		\node[above] at ({\d + \m/2}, 6) {$n$};
		\node[below] at ({6 -\d/2}, 0) {$m-n$};
		\node[below] at ({6 - \d - \m/2}, 0) {$n$};
		\node[above] at (\d/2, 6) {$m-n$};
		\node[left] at (0,\m/2) {$n$};
		\node[left] at (0,{\d/2 + \m}) {$m-n$};
		\node[right] at (6,\d/2) {$m-n$};
		\node[right] at (6,{\m/2 + \d}) {$n$};
		\node[above] at (3,\m) {$2n-m$};
		\node[below] at (3,6-\m) {$2n-m$};
		\node[left] at (6-\m,3) {$2n-m$};
		\node[right] at (\m,3) {$2n-m$};
	\end{tikzpicture}
\end{document}